# Diabetes Risk Factor Analysis - Week 2 Reproducibility Exercise

## Purpose and Overview

This notebook provides an exploratory statistical analysis of a diabetes
risk-factor dataset, which was collected from a cohort of female patients. It was
originally developed for internal review but has now been reorganized into a
reproducible workflow so that a data science team at another clinical site has the
ability to execute it independently and obtain the same results.

With this notebook, a data science team will be able to:

- Load and validate the dataset.
- Summarize descriptive statistics.
- Visualize individual variables and the relationship between variables.
- Test whether two clinical measurements (glucose and BMI) differ significantly
  between patients with and without diabetes.

**How to use this notebook:** run every cell from top to bottom, in order, or
alternatively "Run All." A description of each step is provided as a markdown
cell in each step, along with instructions on how to interpret each output.

## Dataset Source and Provenance

**Dataset:** Pima Indians Diabetes Dataset (provided for this course as
`Example_Dataset_Diabetes.csv`).

**Original source:** National Institute of Diabetes and Digestive and Kidney
Diseases (NIDDK), distributed via the UCI Machine Learning Repository / Kaggle.
The version used in this notebook has been provided by the course as a teaching
dataset.

**Population:** 768 female patients of Pima Indian heritage, age 21 and older.

**Columns:**

| Column | Description |
|---|---|
| Pregnancies | Number of times pregnant |
| Glucose | Plasma glucose concentration (2-hour oral glucose tolerance test), mg/dL |
| D_BP | Diastolic blood pressure, mm Hg |
| Skin_Thickness | Triceps skinfold thickness, mm |
| Insulin | 2-hour serum insulin, mu U/mL |
| BMI | Body mass index, kg/m² |
| Pedigree | Diabetes pedigree function (a score estimating diabetes likelihood from family history) |
| Age | Age in years |
| Outcome | Target variable: 1 = diabetes, 0 = no diabetes |

**Known data-quality caveat:** in this dataset, missing measurements for Glucose,
D_BP, Skin_Thickness, Insulin, and BMI appear to have been encoded as 0 rather
than as a true missing value (a value of 0 mm Hg blood pressure or 0 mg/dL glucose
is not physiologically possible in a living patient). This notebook documents this
pattern in the Data Validation and Data Preparation sections.

**Redistribution:** this dataset is included directly in this project's GitHub
repository (`Example_Dataset_Diabetes.csv`).

## Required Software and Libraries

This analysis uses the following Python libraries:

| Library | Purpose | Version |
|---|---|---|
| pandas | Loading and manipulating tabular data | 2.x |
| numpy | Numerical operations | 1.2x |
| scipy | Statistical tests (t-test, Mann-Whitney U) | 1.1x |
| matplotlib | Plotting | 3.x |
| seaborn | Statistical visualizations (heatmap) | 0.13.x |

If you are running this in Google Colab, all of these libraries are pre-installed,
so installing them is not necessary. If you are running this outside Colab,
uncomment and run the `pip install` line in the Setup and Environment cell once,
before running the rest of the notebook.

A fixed random seed is set in that same cell and reused everywhere randomness
occurs, so that every execution of this notebook will produce the same results.

## Setup / Installation Instructions

**Option A: Google Colab (recommended, no setup required):**
Open the notebook using the Colab link below. All required libraries are
pre-installed in Colab.

**Option B: Local Jupyter / VS Code:**
1. Clone this repository:
   ```
   git clone https://github.com/katerina-smith/Programming_for_HDS_Week2.git
   cd Programming_for_HDS
   ```
2. Launch Jupyter (or open the folder in VS Code) and open the ipynb file.

## How to Execute the Notebook

The cell loading the dataset uses a two-step approach so it works for any user,
in any environment, without needing to edit the code:
1. If the CSV file is present in the current working directory, it is loaded
   directly.
2. If it is not present in the current working directory, it is retrieved
   directly from this project's GitHub repository.

Run every cell from top to bottom, in order:
- In Colab: **Runtime → Restart session and run all**.
- In Jupyter: **Kernel → Restart Kernel and Run All Cells**.

No manual inputs, file uploads, or configuration changes are required.

## Expected Outputs

Running the notebook top-to-bottom produces:

- A dataset validation summary confirming the expected shape, columns, and
  numeric format.
- Descriptive statistics for all numeric variables.
- Visualizations of individual variables and of the relationships between
  variables, including how a patient characteristic relates to diabetes status.
- A Spearman correlation matrix and heatmap.
- Results of two statistical tests comparing patients with and without diabetes:
  an independent-samples t-test (glucose) and a Mann-Whitney U test (BMI).
- A reproducibility check: a random sample of patients and a 
  bootstrap-resampled estimate.

## Dataset Requirements

- `Example_Dataset_Diabetes.csv` must be present in the same folder as the
  notebook (included in this repository), or it will be retrieved automatically
  from this project's GitHub repository.
- No other files are required.

## Assumptions and Limitations

- **Missing data encoded as zero.** Glucose, D_BP, Skin_Thickness, Insulin, and
  BMI all contain physiologically implausible zero values that likely represent
  missingness. This notebook does not impute or remove them; all descriptive
  statistics, tests, and correlations conducted include these zero values. The
  purpose of this exercise is to demonstrate a transparent, reproducible workflow
  structure while still preserving the notebook's original analytical results.
- **Cross-sectional, observational data.** All relationships reported here are
  associations and do not establish causality.
- **Single population.** This dataset reflects one specific patient population.

## Computational Environment

A data scientist at another clinical site should be able to run this notebook by
selecting **Runtime → Restart session and run all** in Colab, using only the
dataset included in this repository. The figures and statistics obtained should
match the results reported in the notebook.

##AI Use

Anthropic. Claude. Used to improve documentation, give examples of optimal 
reproducible workflow, and troubleshoot/optimize code. Also used to generate
 examples of README documentation in GitHub.
